% !TEX root = /home/eb/.vscode/extensions/qftrules.programme-colle/tmp/Exercice.tex

%%%%%%%%%%%%%%%%%%%%%%%%%%%%%%%%%%%
\begin{exo}[PC][1][oral]{Filtrage}
  On considère un signal périodique de fréquence fondamentale $f_0=100 \mathrm{~Hz}$, de valeur moyenne $2 \mathrm{~V}$. II comporte deux harmoniques: pour $n=1$, une amplitude de $0,5 \mathrm{~V}$ et pour $n=3$ une amplitude de $0,2 \mathrm{~V}$. Les autres harmoniques sont nulle
  \begin{quest}
\item Donner l'expression de $e(t)$.
\item Tracer son spectre en amplitude.
\item On fait passer ce signal dans un filtre dont le diagramme de Bode est le suivant :
\end{quest}
\end{exo}
%%%%%%%%%%%%%%%%%%%%%%%%%%%%%%%%%%%

%%%%%%%%%%%%%%%%%%%%%%%%%%%%%%%%%%%
\begin{exo}[PC][1][oral]{Supraconduction}
  On considère le modèle de supraconduction tel que $j=\gamma E$ et $j s=-A / \mu 0 \lambda 2$ lié à la supraconduction (A le potentiel vecteur).
  \begin{quest}
    \item Donner l'équation aux dérivées partielles vérifiée par B
    \item Donner la relation de dispersion et commenter suivant le cas où $\gamma=0$ (faire apparaître une pulsation de coupure) ou $\gamma \neq 0$
\end{quest}
\end{exo}
%%%%%%%%%%%%%%%%%%%%%%%%%%%%%%%%%%%

%%%%%%%%%%%%%%%%%%%%%%%%%%%%%%%%%%%
\begin{exo}[PC][1][oral]{Jonction entre semi-conducteurs}
  On accole deux matériaux semi-conducteurs. Avec la diffusion, des électrons traversent les matériaux. Cela provoque la création d'une distribution volumique de charge notée $\rho_{c}$.

  L'étude sera unidimensionnelle, la permittivité diélectrique du vide sera remplacée par celle du matériau $\varepsilon$ (identique pour les deux) dans les équations de l'électrostatique.

  On donne la figure suivante :

  \begin{center}
  \includegraphics[width=0.4\textwidth]{2023_03_23_e57a46732bffc3502578g-1}
  \end{center}

\begin{quest}
    \item La distribution de charge étant localement neutre, trouver une relation entre $\rho_{1}, \rho_{2}, L_{1}$ et $L_{2}$.

  En déduire que :

  $$
  L_1=\frac{\rho_{2}}{\rho_{2}-\rho_{1}} \ell \text { et } L_2=\frac{\rho_{1}}{\rho_{1}-\rho_{2}} \ell \text { avec } \ell=L_1+L_2
  $$

    \item Donner l'équation de Maxwell-Gauss.
  Déterminer le champ $\E= E(x) \ex$ dans tout l'espace sachant que $\E=\Vec{0}$ quand $x \rightarrow \pm \infty$ et que le champ est continu entre $-L_1$ et $L_2$, .
  Représenter $E(x)$ sur un graphe.

    \item Déterminer le potentiel associé. On prendra $V(0)=0$ pour l’origine des potentiels.

    \item Avec le système précédent, on fabrique un composant dont la caractéristique est $I\left(U\right)$ tel que :

  $$
  I=\mathrm{I}_{\mathrm{S}}\left[\exp \left(\frac{U}{V_{0}}\right)-1\right] \text { avec } \mathrm{I}_{\mathrm{S}}>0 \text { positif. }
  $$

  Quel composant possède une telle caractéristique? Donner une application d'un tel composant.
\end{quest}

\end{exo}
%%%%%%%%%%%%%%%%%%%%%%%%%%%%%%%%%%%

%%%%%%%%%%%%%%%%%%%%%%%%%%%%%%%%%%%
\begin{exo}[PC][1][colle]{Détection synchrone}

  % On note $\rho$ la masse volumique de l'eau.
  %
  % \begin{enumerate}
  % 	\item On donne la loi de l'hydrostatique $\grad P = \rho \vec g$. Déterminer la différence de pression entre deux points dans un liquide au repos en fonction de leur différence d'altitude.
  % 	\item On prend un tonneau de section $S$ et de hauteur  $h$ rempli d'eau. Déterminer la pression de l'eau seule\footnote{c'est-à -dire sans prendre en compte la pression de l'air} $P$ d'un point au fond du tonneau.
  % 	\item On ferme le tonneau et on rajoute un tuyau de section $s$ et de hauteur  $H$. Déterminer la pression de l'eau seule $P'$ sur un point au fond du tonneau.
  % 	\item Pour fixer les idées, calculer $P'/P$ et $m / M$ le rapport des masses d'eau tonneau--tuyau. On prend $H = 3h$ et $s = S / 100$. Manifestez votre surprise à l'égard du résultat.
  % 	\item Montrer la relation donnée en (1) en étudiant un volume de liquide de base $\dx\dy$ situé entre  $z$ et  $z+\dz$.
  % \end{enumerate}
  %
  % \begin{figure}[h]
  % 	\begin{center}
  %
  % 	\begin{tabular}{cc}
  % \begin{tikzpicture}[scale=0.5]
  % 	\draw (0, 8) -- (0, 0) -- (6, 0) -- (6, 8);
  % 	\draw[<->] (7, 0) -- (7, 8) node[midway, right]{$h$};
  % \end{tikzpicture}
  % &
  % \begin{tikzpicture}[scale=0.5]
  % 	\draw (2.5, 11) -- (2.5, 8) -- (0, 8) -- (0, 0) -- (6, 0) -- (6, 8) -- (3.5, 8) -- (3.5, 11);
  % 	\draw[<->] (7, 0) -- (7, 8) node[midway, right]{$h$};
  % 	\draw[<->] (7, 8) -- (7, 11) node[midway, right]{$H$};
  % \end{tikzpicture}
  % 	\end{tabular}
  % 	\end{center}
  % \end{figure}

  % \section{CCINP 2022 Physique - Circuits électriques}
  % \begin{circuitikz}[european]
  % 	\draw (0, 0) node[ground];
  % 	\draw[-latex] (0, 0) -- (0, 2)node[top, left]{$\underline{V_e}$} ;
  % 	\draw (0, 2) to[short, i^=$i$] (0.5, 2) to[R, l={$R$}] (2, 2);
  % 	\draw (2, 0) node[ground] to[C, l={$C$}] (2, 2);
  % 	\draw (2, 2) to[short, i^={$i_s = 0$}] (3, 2);
  % 	\draw (3, 0) node[ground];
  % 	\draw[-latex] (3, 0) -- (3, 2)node[top, right]{$\underline{V_s}$} ;
  % \end{circuitikz}
  %
  %
  % \begin{circuitikz}[european]
  % 	\draw (0, 0) node[ground];
  % 	\draw[-latex] (0, 0) -- (0, 2)node[top, left]{$\underline{V_e}$} ;
  % 	\draw (0, 2) to[short, i^=$i$] (0.5, 2) to[C, l={$C'$}] (2, 2);
  % 	\draw (2, 0) node[ground] to[R, l={$R'$}] (2, 2);
  % 	\draw (2, 2) to[short, i^={$i_s = 0$}] (3, 2);
  % 	\draw (3, 0) node[ground];
  % 	\draw[-latex] (3, 0) -- (3, 2)node[top, right]{$\underline{V_s}$} ;
  % \end{circuitikz}
  %
  % 		Montage \textcircled{\small{1}}  Montage \textcircled{\small{2}}
  %
  % On note $\omega_0 = \frac{1}{RC}$ et $\omega_0' = \frac{1}{R'C'}$.
  %
  % \begin{enumerate}
  % 	\item Calculer $\underline{H}(j\omega)$ pour les deux circuits
  % 	\item Identifier les montages 1 et 2 parmi les réponses en gain suivantes. On note $x = \omega / \omega_0$ pour le montage 1 et $x = \omega / \omega_0'$ pour le montage 2.
  %
  % 		\begin{tabular}{cc}
  % 			\gain{20}{0} & \gain{0}{-20} \\
  % 			$G_a$ & $G_b$ \\
  % 			\gain{20}{-20} & \gain{40}{-40} \\
  % 			$G_c$ &  $G_d$
  % 		\end{tabular}
  %
  % 	\item On effectue le montage suivant. On a $e_2 = E \cos(\omega t + \varphi)$.
  %
  % 	\begin{center}
  % 		\begin{circuitikz}
  % 			\draw (0, 0)
  % 				to[vsourcesin, v^>={$E \cos(\omega t)$}] ++(0, 3)
  % 				to ++(3, 0) ;
  % 			\draw (0, 2) to ++(0, 0.5) to[generic, name=deph] ++(2, 0) to ++(1, 0);
  % 			\node[scale=0.6] at (deph){déphasage};
  % 			\draw[-latex] (2, 1) node[ground] to (2, 3) node[above]{$e_1$};
  % 			\draw[-latex] (2.5, 1) node[ground] to (2.5, 2.5) node[above]{$e_2$};
  % 			\draw (3, 2.25) rectangle node{$\times$} (4.5, 3.25);
  % 			\draw (4.5, 2.75) node[above right]{$ke_1e_2$} -- ++(1, 0)
  % 				to[generic, name=montage] ++(2.5, 0);
  % 			\node[scale=0.6] at (montage){1 ou 2};
  % 			\draw[-latex] (5, 1) node[ground] to ++(0, 1.75);
  % 			\draw (5, 1.375) node[right]{$v'_e(t)$};
  % 			\draw[-latex] (7.5, 1) node[ground] to ++(0, 1.75);
  % 			\draw (7.5, 1.375) node[right]{$v'_s(t)$};
  % 			\draw (0, 0) -- (8, 0);
  % 		\end{circuitikz}
  % 	\end{center}
  %
  % 	\begin{enumerate}
  % 		\item Déterminer $v_e' = ke_1e_2$. Effectuer une analyse fréquentielle.
  % 		\item Exprimer $v_s'$ selon $\omega$, dans le cas où l'on a choisi le montage 1 et dans le cas où l'on a choisi le montage 2.
  % 	\end{enumerate}
  % \end{enumerate}
\end{exo}
%%%%%%%%%%%%%%%%%%%%%%%%%%%%%%%%%%%
